\documentclass[aps,prl,reprint,superscriptaddress,nofootinbib]{revtex4-2}
\usepackage{amsmath,amssymb,amsfonts,mathtools,bm}
\usepackage{graphicx}
\usepackage{booktabs,array}
\usepackage{orcidlink}
\usepackage{xcolor}
\usepackage{hyperref}
\hypersetup{colorlinks=true,linkcolor=blue!55!black,citecolor=blue!55!black,urlcolor=blue!55!black}
\graphicspath{{./}{/mnt/data/}{/mnt/data/ft1/theorem1_feasibility_figures/}{/mnt/data/ft2/}}
\usepackage{cancel}
\usetikzlibrary{arrows.meta,positioning,calc,decorations.pathmorphing,shapes.geometric,fit,backgrounds}

\newtheorem{theorem}{Theorem}
\newtheorem{lemma}{Lemma}
\newtheorem{proof}{Proof}
\newtheorem{proposition}{Proposition}
\newtheorem{corollary}{Corollary}
\newtheorem{definition}{Definition}

\newcommand{\Del}{\mathrm{Del}}
\newcommand{\sym}{\mathrm{sym}}
\newcommand{\Tr}{\operatorname{Tr}}
\newcommand{\Bin}{\mathrm{Bin}}
\newcommand{\Hyp}{\mathrm{Hyp}}
\newcommand{\E}{\mathbb{E}}
\newcommand{\id}{\mathrm{id}}
\newcommand{\ket}[1]{\lvert #1\rangle}
\newcommand{\bra}[1]{\langle #1\rvert}
\newcommand{\braket}[1]{\langle #1\rangle}
\newcommand{\cH}{\mathcal H}
\newcommand{\cN}{\mathcal N}
\newcommand{\cB}{\mathcal B}
\newcommand{\cC}{\mathcal C}
\newcommand{\cI}{\mathcal I}
\newcommand{\cP}{\mathcal P}
\newcommand{\sgn}{\operatorname{sgn}}
\newcommand{\diag}{\operatorname{diag}}
\newcommand{\h}{\mathrm{h.c.}}
\newcommand{\safefig}[3]{\IfFileExists{#1}{\includegraphics[width=#2]{#1}}{\fbox{\parbox{#2}{\centering #3}}}}
\newcommand{\pair}{\mathrm{pair}}
\newcommand{\norm}[1]{\lVert #1\rVert}

\begin{document}
\onecolumngrid
\paragraph*{Group-theoretic interpretation.---}
The binomial PI code can be viewed as a lattice code inside the weight
lattice of the symmetric spin representation.  The symmetric subspace
$\mathcal H_{n,\mathrm{sym}}=\mathrm{Sym}^n(\mathbb C^2)$ is the
spin-$j=n/2$ irreducible representation of the collective $SU(2)$ action,
and the Dicke states are weight vectors for the $U(1)$ subgroup generated by
the excitation-number operator $N=\sum_{i=1}^n |1\rangle\!\langle 1|_i$,
$N\ket{D_k^n}=k\ket{D_k^n}$. Thus $e^{i\theta N}\ket{D_k^n}=e^{i\theta k}\ket{D_k^n}$. The binomial code occupies only the sublattice $k=g\ell$ of this weight
lattice.  The logical label is the character of the coarse parity group
$\mathbb Z_2$, namely $\ell\mapsto (-1)^\ell$: even $\ell$ sectors form
$\ket{0_L}$ and odd $\ell$ sectors form $\ket{1_L}$.  Equivalently, on the
undeleted code space the finite rotation $e^{i\pi N/g}$ acts as logical
$Z$, since $e^{i\pi N/g}\ket{D_{g\ell}^n}=(-1)^\ell\ket{D_{g\ell}^n}$.
Deletion by a branch with $a$ lost excitations maps
$\ket{D_{g\ell}^n}$ to $\ket{D_{g\ell-a}^{n-r}}$.  Since
$a\le r\le t=g-1$, the residue modulo $g$ identifies the branch:
$g\ell-a\equiv -a \pmod g$.  Hence the deletion branch is a charge sector
for the cyclic group $C_g=\langle e^{2\pi iN/g}\rangle$.  The observable
$Z_g$ is the corresponding syndrome-corrected parity character: for
$q=g\ell-a$ it reconstructs $L(q)=\ell$ and returns $(-1)^\ell$.  Thus
$Z_g$ reads the logical $Z$ character perfectly in every $C_g$ syndrome
sector.  By contrast, $X_g$ is an intertwiner between the two
$\mathbb Z_2$ parity sectors inside a fixed $C_g$ charge sector, pairing
weights separated by $g$, $\ket{D_{g\ell-a}^{n-r}}
\longleftrightarrow
\ket{D_{g(\ell+1)-a}^{n-r}}$.
In this language, deletion reveals the $C_g$ syndrome charge but not the
logical $\mathbb Z_2$ character.  The equality
$c_{0,a}^{(r)}=c_{1,a}^{(r)}$ states precisely that the syndrome has equal
weight on the two logical parity sectors.  Therefore each fixed syndrome
branch remains a balanced logical Bell pair, while averaging over syndrome
branches preserves the perfect $Z_gZ_g$ correlation and only reduces the
$X_gX_g$ correlation through the branch coherences $\mu_a^{(r)}$.
%
If the logical state were encoded directly in the "sign of magnetization" (e.g., separating the code space into blocks of high vs. low excitation numbers), the deletion channel would leak information about the logical state. This is because the probability of losing $a$ excitations depends heavily on how many excitations ($gl$) you started with.  
\paragraph{Two finite examples: deletion syndromes and logical parity}

We now illustrate the syndrome-resolved mechanism in two concrete finite
instances.  The first example is deliberately small, so that all branch
states can be written explicitly.  It satisfies the structural conditions
$M$ odd and $M>t$, but not the technical condition $M\ge24$ used later for
the quantitative Hellinger bound.  The second example satisfies all
hypotheses of Theorem~1.

\paragraph*{Example 1: $t=1$, $g=2$, $M=3$, $n=6$.---}
The code is supported on Dicke weights $k=2\ell$, with
$\ell=0,1,2,3$.  Thus
\begin{align}
\ket{0_L}=
\frac{1}{2}\ket{D_0^6}
+
\frac{\sqrt3}{2}\ket{D_4^6},
\qquad
\ket{1_L}=
\frac{\sqrt3}{2}\ket{D_2^6}
+
\frac{1}{2}\ket{D_6^6}.
\end{align}
Consider one deletion, $r=1$.  There are two syndrome branches:
$a=0$, no excitation deleted, and $a=1$, one excitation deleted.  The branch
weights are
\begin{equation}
Q_{\ell,0}^{(1)}
=
\frac{\binom{2\ell}{0}\binom{6-2\ell}{1}}{\binom61}
=
1-\frac{\ell}{3},
\qquad
Q_{\ell,1}^{(1)}
=
\frac{\binom{2\ell}{1}\binom{6-2\ell}{0}}{\binom61}
=
\frac{\ell}{3}
\qquad\Longrightarrow\quad
\begin{array}{c|cccc}
\ell & 0 & 1 & 2 & 3 \\ \hline
Q_{\ell,0}^{(1)} & 1 & 2/3 & 1/3 & 0 \\
Q_{\ell,1}^{(1)} & 0 & 1/3 & 2/3 & 1
\end{array}.
\end{equation}
The squared norms of the two logical sectors in each branch are equal:
\begin{align}
c_{0,0}^{(1)}
=
\frac14\left[
\binom30 Q_{0,0}^{(1)}
+
\binom32 Q_{2,0}^{(1)}
\right]
=
\frac14(1+1)
=
\frac12,
\qquad
c_{1,0}^{(1)}
=
\frac14\left[
\binom31 Q_{1,0}^{(1)}
+
\binom33 Q_{3,0}^{(1)}
\right]
=
\frac14(2+0)
=
\frac12,
\end{align}
and similarly 
$c_{0,1}^{(1)}=c_{1,1}^{(1)}=\frac12$.
Thus the deletion syndrome $a$ is observed with probability $1/2$ in each
logical sector; it does not reveal the logical parity of $\ell$.
%
The unnormalized branch states are
\begin{align}
\ket{0_0^{(1)}}=
\frac{\ket{D_0^5}+\ket{D_4^5}}2,
\quad
\ket{0_1^{(1)}}&=
\frac{\ket{D_3^5}}{\sqrt2},
\qquad
\ket{1_0^{(1)}}=
\frac{\ket{D_2^5}}{\sqrt2},
\quad
\ket{1_1^{(1)}}=
\frac{\ket{D_1^5}+\ket{D_5^5}}2.
\end{align}
After normalization,
\begin{align}
\ket{\widetilde0_0^{(1)}}=
\frac{\ket{D_0^5}+\ket{D_4^5}}{\sqrt2},
\quad
\ket{\widetilde0_1^{(1)}}=\ket{D_3^5},
\qquad
\ket{\widetilde1_0^{(1)}}=\ket{D_2^5},
\quad
\ket{\widetilde1_1^{(1)}}=
\frac{\ket{D_1^5}+\ket{D_5^5}}{\sqrt2}.
\end{align}

This example makes the residue-class structure explicit.  In branch $a=0$,
the surviving Dicke weights are even, $q=2\ell$, whereas in branch $a=1$
they are odd, $q=2\ell-1$.  Since $a<g$, the residue of $q$ modulo $g$
identifies the deletion syndrome.  At the same time,
$L(q)=\frac{q+a(q)}{2}=\ell$,
so $Z_2$ still reads the logical parity:
$Z_2\ket{D_{2\ell-a}^{5}}
=
(-1)^\ell\ket{D_{2\ell-a}^{5}}$.
The deletion branch changes the residue class modulo $2$, but it does not
change the coarse parity label $\ell\bmod2$.
%
The branch coherences are 
$\mu_0^{(1)}=\mu_1^{(1)}=\frac{1}{\sqrt2}$,
and $\mu_1=\frac12\mu_0^{(1)}+\frac12\mu_1^{(1)}
=
\frac{1}{\sqrt2}$.
Therefore, if both parties suffer one deletion,
$S_{1,1}
=
\sqrt2\left(1+\mu_1^2\right)
=
\sqrt2\left(1+\frac12\right)
=
\frac{3\sqrt2}{2}>2$.
Thus the averaged post-deletion state is not a single pure Bell state, but
for the observables $Z_2$ and $X_2$ it behaves as a Bell pair with perfect
$Z_2Z_2$ correlation and reduced $X_2X_2$ visibility $1/2$.

\paragraph*{Example 2: $t=2$, $g=3$, $M=25$, $n=75$.---}
We now consider a theorem-compatible instance.  Here $M$ is odd, $M\ge24$,
and $M>t$.  The code is supported on Dicke weights
$k=3\ell$, for $\ell=0,1,\ldots,25$.
The logical label is the parity of $\ell$, while the deletion syndrome is
the residue class of the surviving Dicke weight modulo $3$.
%
Consider the worst correctable fixed deletion number, $r=2$.  There are
three deletion syndromes, $a=0,1,2$.  The corresponding branch weights are
\begin{align}
Q_{\ell,0}^{(2)}
=
\frac{\binom{3\ell}{0}\binom{75-3\ell}{2}}{\binom{75}{2}}
=
\frac{(75-3\ell)(74-3\ell)}{5550},
\quad
Q_{\ell,1}^{(2)}
=
\frac{\binom{3\ell}{1}\binom{75-3\ell}{1}}{\binom{75}{2}}
=
\frac{3\ell(25-\ell)}{925},
\quad
Q_{\ell,2}^{(2)}
=
\frac{\binom{3\ell}{2}\binom{75-3\ell}{0}}{\binom{75}{2}}
=
\frac{(3\ell)(3\ell-1)}{5550}.\nonumber
\end{align}
Each $Q_{\ell,a}^{(2)}$ is a polynomial in $\ell$ of degree at most $2$.
Since $2<M=25$, the finite-difference identity gives
$c_{0,a}^{(2)}=c_{1,a}^{(2)}$
for $(a=0,1,2)$.
In this particular example, one finds 
$c_0^{(2)}=c_2^{(2)}=\frac{19}{74}$,
and $c_1^{(2)}=\frac{18}{37}$.
Again, the syndrome distribution is independent of the logical parity:
deletion distinguishes the charge sector $a\bmod3$, but not whether
$\ell$ is even or odd.
%
The group-theoretic structure is now especially transparent.  A component
$\ket{D_{3\ell}^{75}}$ is mapped in syndrome branch $a$ to
$\ket{D_{3\ell-a}^{73}}$.
Since $a=0,1,2<3$, the residue of the surviving Dicke weight determines
the syndrome:
$3\ell-a\equiv -a \pmod 3$.
Thus the three branches lie in the three character sectors of the cyclic
group $C_3$.  Within each such sector, the logical observable $Z_3$ reads
the remaining $\mathbb Z_2$ character:
$Z_3\ket{D_{3\ell-a}^{73}}
=
(-1)^\ell\ket{D_{3\ell-a}^{73}}$.
Meanwhile $X_3$ preserves the syndrome charge and pairs neighboring coarse
weights, $\ket{D_{3\ell-a}^{73}}
\longleftrightarrow
\ket{D_{3(\ell+1)-a}^{73}}$.
Thus $Z_3$ is a syndrome-corrected parity character, while $X_3$ is an intertwiner between the even- and odd-$\ell$ sectors inside each fixed $C_3$ syndrome sector.
For $r=2$, the branch coherences are
\begin{equation}
    \mu_k^{(2)}=
\frac{1}{c_k^{(2)}2^{24}}
\sum_{\ell\ \mathrm{even}}
\sqrt{
\binom{25}{\ell}\binom{25}{\ell+1}
Q_{\ell,k}^{(2)}Q_{\ell+1,k}^{(2)}
},\qquad k=0,1,2.
\end{equation}
Numerically, 
$\mu_0^{(2)}\simeq 0.97855$3,
$\mu_1^{(2)}\simeq 0.978457$,
$\mu_2^{(2)}\simeq 0.978553$.
Averaging over the three deletion syndromes gives 
$\mu_2
=
\sum_{a=0}^2 c_a^{(2)}\mu_a^{(2)}
\simeq
0.978507$.
Therefore, for fixed deletion numbers $r_A=r_B=2$,
$S_{2,2}
=
\sqrt2\left(1+\mu_2^2\right)
\simeq
2.76829$.
\end{document}